\scanpage{51}
\mspage{\deleted{\unsure{6}}7}
% archive sheet number: 50
d'une structure déjà orientée si $a$ est une arête libre (correspondant à un $C\in\pi_{0}(\partial X)$), ou qu'on dispose \inserted{de} l'une des deux orientations $\omega(a)$ --- dans le cas où $a$ est une arête ordinaire, indique que l'on tient en mains, pour chaque $a\in A$, un polygone combinatoire $P_a$, dont l'ens des sommets est $R_a$, et l'ens des arêtes peut se décrire \deleted{\unsure{comme}} de façon canonique en termes de $R_a,\omega_a$ et $\omega_a\longrightarrow\mathrm{Circ}(R_a)$. Il peut être plus commode de le \unsure{donner|sortir} d'emblée dans la structure, où la famille des $R_a$, serait remplacée par un \deleted{\unsure{qui serait donc un ensemble}} \uline{contour combinatoire}, i.e un ensemble de repères $\underline{\vec R}$ (plutôt ``repères de rigidification''), avec une opération de
\[
\underline{\unsure{D|\mathbb D}}_{\infty}
=\{\sigma_{0},\sigma_{1}\mid\sigma_{0}^{2}=\sigma_{1}^{2}=1\}
\]
dessus, et une application
\diagram{
\begin{tikzcd}[ampersand replacement=\&]
\underline{\vec R}\arrow[r]\&\vec A
\end{tikzcd}
}
compatible avec l'hom
\diagram{
\begin{tikzcd}[ampersand replacement=\&]
\underline D_{\infty}\arrow[r,"\unsure{\varepsilon|\epsilon}"]\&\{\pm 1\}
\end{tikzcd}
}
sur les groupes d'opérateurs, mais de plus un sous ens
\[
\underline{\vec R}_{\ell}^{+}\subset\underline{\vec R}|\,\vec A_{\ell}
\qquad
\begin{aligned}
&(\vec A_{\ell}=\text{arêtes libres,}\\
&\quad\text{correspondant aux}\\
&\quad\text{cercles de \unsure{recollement|recollage}}\\
&\quad\text{i.e. les }\in\pi_0(\deleted{\unsure{\Sigma}}\partial X))
\end{aligned}
\]
définissant une orientation des contours \unsure{correspondant|correspondants}.
